\section{Solution Overview}

FalconVQA handles a video in five steps: split it into chunks, analyze each chunk,
index the results, roll them up into video-level summaries, then search or ask
questions. Most steps have more than one setting, because footage varies too much for
one fixed choice to work everywhere.

\optfigure{video-mind-arch.png}{FalconVQA system architecture.}{1.0}

\subsection{Splitting the video}

We cut the video where something changes. Four detectors look for that: a change of
speaker, a stretch of silence, a hard cut, and a shift in what the frame looks like.
Their scores are combined and the strongest peaks become the cut points.

If a signal isn't there at all, its share is passed to the ones that did fire. This
matters more than it sounds: on a fixed camera there are no cuts and often no speech,
and without that adjustment the whole clip comes back as one chunk. There are two
other modes for when this isn't wanted --- supply your own weighting, or just cut
every $N$ seconds.

The chunking settings are stored alongside the results, so the same video can be
indexed coarsely and finely at the same time, and a search picks whichever it needs.

\subsection{Analysing each chunk}

You pick which analyzers to run when you upload, because they do not cost the same.
There is one for describing the scene, one for people, one for objects, one for
on-screen text, and one for speech.

The vision model is the expensive part, so cheap detectors run first and act as a
filter. If YOLO sees nothing in a frame and EasyOCR finds no text, that frame is never
sent. Frames are also deduplicated: we keep a frame only if it looks different enough
from the last one we kept. The threshold adapts, with a ceiling --- a fixed value kept
just one frame on a static camera, and a purely relative one started picking up
compression noise. Frames are read in order rather than seeked to, which measured
about $22\times$ faster.

\subsection{Searching and asking}

Each chunk is stored as several separate vectors instead of one: the full record, the
scene description, the people, the actions, the objects. So a search about what
someone looked like is matched against a short description of a person, rather than
competing with everything else said about that chunk. Searches can be filtered by
time, objects, tags, people and so on, and can return more or less detail depending on
who is reading them.

Questions take a different path. Rather than searching chunks and hoping, the question
is matched to the video-level summaries that can actually answer it: the entity
narratives for a question about a person, novelty for ``anything odd here'', the
statistics for ``how busy was it''.

Either way the operator types a description and gets back ranked moments with
\texttt{mm:ss} timecodes. Clicking one plays the video from that point.
